% Korean Noether paper 01 — mechanically assembled inherited production body.
% Reopened against ED0004 by the current producer; independent checking remains external.

% inherited: noether_paper01_ko_translation_001_20260804\ko\Noether_Paper01_Korean_U01_translation_draft_v001.tex
% sha256: 48961F41A3C178968A5D2157F6FD5E756DAC7817555CAD07208C61E5A6643BE7
\begin{center}
{\Large\bfseries 1. 삼원 4차 형식의 형식계 구성에 관하여.}\par
\vspace{1.2em}
에미 뇌터(Emmy Noether).\par
\vspace{0.8em}
(저자의 학위논문에서 발췌.)\par
\vspace{1.2em}
Sitz. Ber. d. Physikal.-mediz. Sozietät in Erlangen 39 (1907), 176--179쪽
\end{center}

삼원 4차 형식의 형식계는 고르단(Gordan), 마이사노(Maisano), 파스칼(Pascal)의 논문들\srcfn{1)}{\srcspaced{P. Gordan}, Über das volle Formensystem der ternären biquadratischen Form $f=x_1^3x_2+x_2^3x_3+x_3^3x_1$ (Math. Annalen Bd. XVII, S. 217 bis 233. 1880); \srcspaced{G. Maisano}, 1. Sistemi completi dei primi cinque gradi della forma ternaria biquadratica e degl' invarianti, covarianti e contravarianti di sesto grado (Giorn. di Battaglini XIX). 2. Sui covarianti indipendenti di $6^\circ$ grado nei coefficienti della forma biquadratica ternaria (Rend. Circ. Mat. di Palermo I. 1887); \srcspaced{E. Pascal}, Contributo alla teoria della forma ternaria biquadratica e delle sue varie decomposizioni in fattori (Memoria premiata dalla R. Accademia delle scienze fisiche e matematiche di Napoli. 1905).}에서 다루었다. 고르단은 특수한 자기동형 형식
$f=x_1^3x_2+x_2^3x_3+x_3^3x_1$의 54개 구성식으로 이루어진 완전한 형식계를 세웠는데, 그 바탕에는 그가 이원 형식론의 형식계에 사용한 것과 비슷한 원리들이 놓여 있다.

마이사노는 일반 4차 형식에 대하여 계수차수 5까지의 형식들\srcfn{2)}{여기서 ``계수차수(Ordnung)''는 계수에 관한 차수를, ``변수차수(Grad)''는 변수에 관한 차수를 뜻한다.}과 더 높은 계수차수의 몇몇 불변식, 공변식, 반변식을 제시했다. 그가 사용한 방법은 고르단이 \emph{Mathematische Annalen} 제1권에서 삼원 3차 형식에 적용한 방법이다. 파스칼은 마이사노의 결과를 이용하면서 주로 4차 형식이 인수들로 분해되는 문제를 다룬다\srcfn{1)}{두 번째 짧은 논문에서 마이사노는 계수차수 6, 변수차수 6인 세 공변식의 선형종속성을 증명하려 한다. 완전히 갖추어지지 않은 이 증명과 달리 파스칼은 계수차수 6인 세 공변식의 선형독립성과 계수차수 9인 세 불변식의 선형독립성을 증명했다고 본다. 그러나 실제로는 한편의 세 공변식 사이와 다른 한편의 세 불변식 사이에 각각 선형관계가 존재한다는 것이 나의 계산에서 드러났다. 파스칼의 표기법으로 쓴 명시적 공식은 다음과 같다.
\[
0=20\Omega_1+6\Omega_2-3\Omega_3+4f\cdot C_1-12f\cdot C_2-4A\cdot\Delta.
\]
\[
0=4A^3-15A\cdot B+30C-90D+9E.
\]}

\emph{내 연구의 목표는 일반 삼원 4차 형식의 형식계를 세우는 것이다. 여기서는 우선 그 주요 토대만을 제시하고, 이른바 ``상대적으로 완전한 계''를 세운다.}\srcfn{2)}{명칭에 관해서는 \srcspaced{Gordan-Kerschensteiner}, Vorlesungen über Invariantentheorie, Bd. II, S. 227을 보라.}

삼원 형식의 계를 구성하는 기본 발상은 이원 형식론에서와 같다. 첫 번째 상대적으로 완전한 계, 곧 이원 형식론에서 가져온 형식들의 계에서 출발하여 일정한 법칙에 따라 모듈이 갈수록 높은 계들로 나아간다. 이 과정은 어떤 모듈의 계가---그 모듈을 기본형식으로 삼았을 때---유한하고 알려진 것이 되거나, 또는 어떤 모듈이 불변식을 인수로 갖는 형식들로 환원될 수 있을 때까지 계속된다. 첫째 경우에는 상대적으로 완전한 계와 그 모듈의 계에 전이연산(Überschiebung, transvection)을 시행하여 절대적으로 완전한 계를 얻는다. 둘째 경우에는 상대적으로 완전한 계와 절대적으로 완전한 계가 일치한다. 힐베르트(Hilbert)가 형식계의 유한성을 일반적으로 증명했으므로, 이 절차는 반드시 끝나야 한다.

% inherited: noether_paper01_ko_translation_001_20260804\ko\Noether_Paper01_Korean_U02_translation_draft_v001.tex
% sha256: 52C02759CC6D08AA102DA366F7F148A4D148EC1066E2F81DE929CEE43A46DDDF
우리의 경우에는 첫 번째 상대적으로 완전한 계의 모듈 $(abc)$를 다음 모듈들로 환원할 수 있다.
\[
\Delta=(abc)^2a_x^2b_x^2c_x^2\quad\text{및}\quad \nu=(abu)^4.
\]
이로부터 mod $\nu$에 관한 상대적으로 완전한 계를 쉽게 얻는다. 이제 모듈열로 다음 형식들을 택한다.
\[
\nu;\qquad \nu^{(\nu)}=(\nu\nu_1x)^4=s_x^4,
\]
\[
\nu(s)=(ss'u)^4\ldots
\]
그리고 mod $s$에 관한 상대적으로 완전한 계를 구성할 때 나타나는 두 이차형식 $u_\varrho^2$와 $t_x^2$를 추가 모듈로 덧붙인다. 그러면 $s$ 다음에 오는 모듈 $(ss'u)^4$가 모듈 $(\varrho,t)$로 환원 가능함을 보일 수 있다. 그런데 두 이차형식의 동시 형식계는 유한하고 알려져 있으므로, 이로써 앞에서 말한 종결에 도달한다. mod $(\varrho,t)$에 관한 상대적으로 완전한 계로는 331개의 구성식이 얻어진다. ---

이제 사용한 방법에 관하여 몇 가지를 더 적어 둔다.

형식을 생성하는 기본 과정인 \emph{수축 과정(Faltungsprozeß)}은 삼원 형식론에서 다음과 같이 정의할 수 있다.

기호적 곱
\[
s_x^m t_x^n u_\sigma^\mu u_\tau^\nu
\]
에서 인수쌍들을
\[
\begin{array}{@{}c@{\qquad}c@{\qquad}c@{\qquad}c@{\qquad}c@{}}
 & s_xt_x & u_\sigma u_\tau & s_xu_\sigma\ \text{또는}\ s_xu_\tau & t_xu_\tau\ \text{또는}\ t_xu_\sigma\\[0.45em]
\text{각각 다음으로} & (stu) & (\sigma\tau x) & s_\sigma\ \text{또는}\ s_\tau & t_\tau\ \text{또는}\ t_\sigma\\[0.45em]
\text{수축} & \mathrm{I} & \mathrm{II} & \mathrm{III} & \mathrm{IV}
\end{array}
\]
바꾸면, 그 결과로 생기는 형식들은 원래 형식에서 수축을 통해 얻어진 것이다\srcfn{1)}{\srcspaced{Gordan}. Math. Annalen, Bd. XVII, S. 219.}.

이때 언제나 $s_\sigma=0;\ t_\tau=0$으로 둘 수도 있다. 각 수축 사이의 관계에 대해서는 다음 정리가 성립한다.

\emph{수축 I과 II는 기본 수축이다. 수축 III과 IV는 합성 순서와 무관하게 이 기본 수축들로 합성할 수 있다. 다시 말해 주어진 식에서 수축으로 생기는 모든 형식을 만들려면 수축 I과 II만 적용하면 된다.}\srcfn{2)}{$n$과 $\mu$, 또는 각각 $m$과 $\nu$가 동시에 0이 되는 경우에는 이 정리에 예외가 생긴다. 그때 ``형식열''은 대각 성분 하나로 줄어든다.}

% inherited: noether_paper01_ko_translation_001_20260804\ko\Noether_Paper01_Korean_U03_translation_draft_v001.tex
% sha256: ECEE0AB9E9D8C89D6A9B4FBBA63128FBE1990764847E01038AB894EF66C9DF54
``형식열(Formenreihe)''이란 하나의 시작형식과, 그 형식 자체에서 수축을 통해 생겨나는 모든 형식을 함께 이르는 것으로 정의한다. 위 정리에 따르면 형식열 $s_x^m t_x^n u_\sigma^\mu u_\tau^\nu$를 직사각형 도식으로 배열할 수 있다. 도식에서 나아갈 때에는 다음을 얻는다.
\begin{enumerate}
\item 오른쪽으로 한 열 이동하면, 바로 왼쪽의 형식들에서 수축 I로 생기는 모든 형식,
\item 아래쪽으로 한 행 이동하면, 바로 위의 형식들에서 수축 II로 생기는 모든 형식, 따라서 수축으로 생기는 모든 형식.
\end{enumerate}

이 전제 아래에서 환원자에 관한 정리들은 가장 일반적인 형태로 말할 수 있다. 먼저 다음과 같이 정의한다.
\begin{center}
``환원자(Reduzent)란 환원 가능한 형식열이다.''
\end{center}
그러면 다음이 성립한다.

\begin{center}
\emph{어떤 형식열의 시작형식이 환원 가능한 까닭이 그 성분 하나가 환원자와의 수축으로 생겨났기 때문이고, 그 형식열의 끝형식이 바로 그 성분에서 수축으로 생겨났다면, 전체 형식열은 환원 가능하다.}
\end{center}

그 밖의 환원 방법으로는 다음을 들 수 있다.
\begin{enumerate}
\item 이른바 ``이중 환원'', 곧 한 형식을 두 가지 방식으로 환원하여 더 높은 형식들 사이의 관계를 얻는 방법,
\item ``분해되는 형식과의 수축''. 이 방법은 일부는 환원자를 통한 환원의 성격을, 일부는 ``이중 환원''의 성격을 띤다.
\end{enumerate}
